\section{Future Work}\label{sec:future}

The remaining engineering work is partitioned into three named
milestones; each corresponds to an annotated GPG-signed git tag
and to a subsequent version of this report.

\begin{description}
\item[v0.3.0-metadata-hardening (target: v2 of this report).]
Implementation of item~4: format~v4 bump, struct
\texttt{ftrfs\_rs\_event} extended to 40~bytes, Shannon entropy
field with Q16.16 fixed-point lookup table, superblock layout
migration to 13~RS subblocks. The migration is breaking: a v4
kernel refuses to mount a v3 image (per axis~3.3). v2 will report
the recovery delta against the v1 baseline of \cref{tab:iobench}.

\item[v0.4.0-universal-protection (target: v3).]
Stage~4: universal data-block protection. The choice between
\texttt{UNIVERSAL\_INLINE} (parity bytes embedded per block),
\texttt{UNIVERSAL\_SHADOW} (out-of-band parity region), and
\texttt{UNIVERSAL\_EXTENT} (extent-attribute parity) will be made
on the basis of measured throughput, recovery latency, and
integration cost. v3 will additionally include the deferred
user-space CVE remediation.

\item[v0.5.0-security-reviewed (target: v4).]
Stage~5: targeted syzkaller fuzzing campaign with a custom
\texttt{syzlang} grammar for the FTRFS mount and ioctl surfaces;
KASAN-instrumented stress runs; structured adversarial review
against the threat model of \cref{sec:threat}.
\end{description}

Each version of this report is deposited on HAL with the same
\texttt{idHAL} (\texttt{aurelien-desbrieres}); HAL native
versioning preserves all prior versions and timestamps each new
submission, providing cumulative anteriority. Each version is
mirrored to Zenodo via a GitHub release, yielding a perennial DOI.
The source LaTeX, bibliography, dataset, and build instructions
for the present report are committed under
\texttt{papers/2026-04-ftrfs-v1/} in the FTRFS source
tree~\cite{ftrfs-repo}.
